\section{Выводы}

Выполнив пятую лабораторную работу, я глубоко изучил теорию и практику суффиксных деревьев, а также применил полученные знания для решения задачи о лексикографически минимальной циклической ротации строки.

В процессе выполнения лабораторной работы я:

\begin{itemize}
    \item Полностью разобрал алгоритм Укконена построения суффиксного дерева за линейное время, включая механизм активного состояния и суффиксные ссылки.

    \item Реализовал эффективный поиск оптимальной ротации через обход дерева с выбором лексикографически минимального переходу на каждом шаге.

    \item Применил классический приём конкатенации $s + s + \$$ для представления всех циклических ротаций как суффиксов в одной структуре.

    \item Провел бенчмарки, демонстрирующие практическое превосходство suffix tree ($O(n)$) над наивным подходом ($O(n^2)$) на строках разной длины.
\end{itemize}

Реализация демонстрирует корректное понимание структур данных на строках, эффективный алгоритм построения суффиксного дерева и применение этой структуры для решения практических задач оптимизации. Полученные навыки применимы к широкому спектру задач обработки текстов, включая поиск наибольшей общей подстроки, определение повторяемых элементов и анализ строковых данных.

\pagebreak
